\chapter*{Notation}
\addcontentsline{toc}{chapter}{Notation}
\label{front:notation}

Symbols are defined at first use. The tables below are the standing vocabulary, grouped so that a later chapter can point back here---and so that a symbol here can point forward to the chapter that uses it.

This list will grow. It is not a course in any of the subjects named.

\begingroup
\renewcommand{\thetable}{N.\arabic{table}}

\section*{Calculus}
\label{not:calculus}

The intended start is a first university course in calculus. Table~\ref{tab:notation-calculus} is the core; derivatives and integrals return in Chapters~\ref{ch:optimization} and~\ref{ch:dynamics}.

\begin{table}[htbp]
\centering
\moicaption{Calculus}{Maps, derivatives, integrals, and asymptotic notation. Later chapters point here when a rate, a gradient, or an integral is the object in play.}
\label{tab:notation-calculus}
\begin{tabular}{@{}ll@{}}
\toprule
Symbol & Meaning \\
\midrule
$\N,\Z,\R,\C$ & Natural, integer, real, and complex numbers \\
$f\colon \R^n\to\R^m$ & A map from an $n$-vector to an $m$-vector \\
$\partial_i f$, $\grad f$ & Partial derivative; gradient \\
$\Hess f$, $Hf$ & Hessian of a scalar $f$ \\
$Df$, $J_f$ & Derivative / Jacobian of a vector-valued map \\
$\dfrac{\mathrm{d}}{\mathrm{d}t}$ & Ordinary derivative in a scalar parameter \\
$\int_{\Omega} f$ & Integral of $f$ over a domain $\Omega$ \\
$o(\cdot)$, $O(\cdot)$ & Little-o and big-O asymptotics \\
$\arg\min$, $\arg\max$ & Minimiser and maximiser of an objective \\
\bottomrule
\end{tabular}
\end{table}

\section*{Matrices and vectors}
\label{not:linear}

Inner products and norms are the language of Chapter~\ref{ch:spaces-and-inner-products}. Spectra and decompositions are taken up in Chapter~\ref{ch:operators-and-spectra}.

\begin{table}[htbp]
\centering
\moicaption{Matrices and vectors}{Linear algebra used as a language for representation, scores, and parameter updates. The geometry of $\langle u,v\rangle$ is Figure~\ref{fig:inner-product} in Chapter~\ref{ch:spaces-and-inner-products}.}
\label{tab:notation-linear}
\begin{tabular}{@{}ll@{}}
\toprule
Symbol & Meaning \\
\midrule
$x\in\R^d$ & A real vector of dimension $d$ \\
$x^\top$, $x_i$ & Transpose; $i$-th coordinate \\
$\langle u,v\rangle$ & Inner product \\
$\lVert x\rVert$, $\lVert x\rVert_p$ & Induced norm; $\ell^p$ norm \\
$A\in\R^{m\times n}$ & A real matrix with $m$ rows and $n$ columns \\
$A^\top$, $A^{-1}$, $A^{+}$ & Transpose, inverse, and Moore--Penrose pseudoinverse \\
$\tr A$, $\det A$, $\mathrm{rank}\,A$ & Trace, determinant, and rank \\
$I_d$, $\diag(v)$ & Identity of size $d$; diagonal matrix from a vector \\
$\lVert A\rVert_2$, $\lVert A\rVert_F$ & Operator (spectral) norm; Frobenius norm \\
$\lambda(A)$, $\sigma(A)$ & Eigenvalues; singular values \\
$uv^\top$, $A\otimes B$ & Outer product; Kronecker product \\
\bottomrule
\end{tabular}
\end{table}

\section*{Probability, measure, and function spaces}
\label{not:probability}

Expectation and conditionals are fixed in Chapter~\ref{ch:probability}. Information quantities return in Chapter~\ref{ch:information}. Hilbert-space language is used whenever a score is an inner product, including later generative chapters.

\begin{table}[htbp]
\centering
\moicaption{Probability, measure, and function spaces}{Random objects, measures, and $L^p$ / Hilbert spaces. Measure-theoretic and functional-analytic names appear only when a later argument needs them.}
\label{tab:notation-probability}
\begin{tabular}{@{}ll@{}}
\toprule
Symbol & Meaning \\
\midrule
$(\Omega,\mathcal{F},\Prob)$ & A probability space \\
$X\sim p$ & A random variable with law $p$ \\
$p(x)$, $p(x\mid y)$ & Density or mass; conditional \\
$\E[X]$, $\Var(X)$, $\Cov(X,Y)$ & Expectation, variance, and covariance \\
$\ind_A$ & Indicator of a set $A$ \\
$\delta_x$ & Dirac measure at $x$ \\
$\mu$, $\dfrac{\mathrm{d}\mu}{\mathrm{d}\nu}$ & A measure; Radon--Nikodym derivative \\
$L^p(\mu)$ & Functions with integrable $p$-th power \\
$\Hilbert$ & A Hilbert space \\
$\langle\cdot,\cdot\rangle_{\Hilbert}$ & Inner product on $\Hilbert$ \\
$T^{*}$ & Adjoint of a linear operator $T$ \\
$\KL(p\,\|\,q)$ & Kullback--Leibler divergence \\
\bottomrule
\end{tabular}
\end{table}

\section*{Integral transforms}
\label{not:transforms}

Transforms turn a function on one domain into a function on another. They will be invoked from dynamics, scores, and generative models---in particular Chapter~\ref{ch:dynamics} and Chapter~\ref{ch:diffusion}---rather than treated as a separate part.

\begin{table}[htbp]
\centering
\moicaption{Integral transforms}{Fourier, Laplace, convolution, and the characteristic function. Each is a change of coordinates, not a new kind of object.}
\label{tab:notation-transforms}
\begin{tabular}{@{}ll@{}}
\toprule
Symbol & Meaning \\
\midrule
$\Fourier f$, $\hat{f}$ & Fourier transform of $f$ \\
$\Fourier^{-1}\hat{f}$ & Inverse Fourier transform \\
$\Lapl\{f\}$ & Laplace transform of $f$ \\
$(f*g)(x)$ & Convolution of $f$ and $g$ \\
$\varphi_X(t)$ & Characteristic function $\E[e^{\mathrm{i}tX}]$ \\
$M_X(t)$ & Moment-generating function $\E[e^{tX}]$, when it exists \\
\bottomrule
\end{tabular}
\end{table}

\section*{Objects in later parts}
\label{not:objects}

These names are not mathematics in themselves. They label the things Part~\ref{part:mathematics-of-ai} and Part~\ref{part:intelligence-and-beyond} manipulate.

\begin{table}[htbp]
\centering
\moicaption{Objects in later parts}{Loss, data, and model class as they appear from Part~\ref{part:mathematics-of-ai} onward. They sit on top of the tables above, not beside them.}
\label{tab:notation-objects}
\begin{tabular}{@{}ll@{}}
\toprule
Symbol & Meaning \\
\midrule
$\loss$ & Loss or objective \\
$\data$ & Data set or empirical measure \\
$\model$ & Model class \\
\bottomrule
\end{tabular}
\end{table}

\endgroup
